\section{Data-collection system}
\label{sec:collection}

The system produces two artefacts per session: a training-ready dataset and
a full-rate diagnostic log. This section walks the two-rate design, the
stamp-on-read rule, the reference-time alignment, the drift budget, the
recorded stream set, and the phase flag and gating.

\paragraph{Two-rate design.} A dataset carries a single frame
rate~\citep{lerobot}, so the training-ready features are recorded at one
fixed rate suited to the policy. Diagnosis, however, needs the control signals at their native, much
higher rate. We therefore keep a separate full-rate side log with
wall-clock stamps that holds every control signal, stored beside the
dataset and never fed to the policy. The dataset trains; the side log
explains.

\paragraph{Stamp on read.} Each stream is timestamped at the instant it is
read from its device, before any decoding or resizing, on a monotonic
clock. Stamping at read rather than at publication means the recorded time
reflects when the sensor was actually sampled, not when the processed
result happened to arrive.

\paragraph{Reference-time alignment.} When the recorder builds a frame it
fixes one reference time and samples every stream at that time. Images take
the nearest captured frame; proprioception and end-effector poses are
interpolated to the reference time, with orientations interpolated along
the shortest rotation. This replaces the naive rule of taking each stream's
latest value, which would place streams captured at different instants into
one frame as though they were simultaneous.

\paragraph{Drift budget.} For every stream and every frame the collector
records the residual between the reference time and that stream's nearest
real sample, and a live monitor shows these residuals during collection.
This matters because unsynchronised cameras cannot be re-phased in
software: their capture instants are fixed by their own clocks. The honest
deliverable is therefore a \emph{measured} drift, not a claim of perfect
synchronisation. We target a per-stream drift well within half a frame
period for the proprioceptive streams; the cross-camera residual is
whatever the hardware imposes, and the monitor makes it visible so an
operator can react to a starved bus. The same monitor also carries the
control prompts and a live progress readout --- the episodes recorded
against the session's target, together with the recorder's current phase, so
a save in flight is visible on the window --- so an operator drives a whole
session from the window rather than the terminal. \unjustified{} the half-frame target
is a working rule of thumb, not a value validated against downstream policy
performance.

\paragraph{Capture budget.} The bus, not only the clock, sets that residual.
Several colour streams share the host's controllers, and an uncompressed
transport does not fit: one uncompressed view at our resolution and rate needs
roughly eighteen megabytes per second, so a few of them exceed what the
controllers carry. The cameras then negotiate a fraction of the rate we ask
for. Nothing reports an error, because each frame that does arrive is valid;
the newest frame is simply already half an inter-arrival gap old, and the
measured residual reflects the capture rate rather than clock skew. We
therefore require a compressed transport on every stream, and a driver queue
short enough to bound staleness yet deep enough not to starve the camera. Both
halves are measured, not assumed. Requesting an uncompressed format held our
wrist cameras to under a third of their rate and left them a mean of
\SI{54}{\milli\second} behind the reference time; compressing them recovered
the full rate and cut that residual to \SI{20}{\milli\second}. A single-frame
queue looks attractive for freshness but halves the achieved rate, because the
driver has nowhere to place the next frame while the reader holds the only
buffer; two buffers sustain the full rate.

Fixing the transport does not exhaust the problem, because a camera's exposure
governs its rate as directly as the bus does: no camera can deliver frames
faster than it exposes them, so an exposure that grows past the frame period
forces the sensor to a slower one. Left automatic, exposure grows exactly when
the scene is dim, and the wrist views are the dim ones by construction --- they
look down at a close workspace that the arm itself shadows. The rate then
depends on the room rather than on the configuration, which is the property that
makes it hard to notice: the same rig measured in daylight and in the evening
gives different answers, and neither reports a fault. On our cameras the
relation is steep. Held to an exposure short enough to fit inside the frame
period, the wrist streams sustain their full rate; one step longer costs a fifth
of it; and the value automatic exposure selected under collection lighting gave
two thirds of the frame period to exposure alone, leaving the stream at little
over half its rate --- which was precisely the rate those streams had been
recording at. We therefore fix the exposure of the streams whose scene is
predictably dark. The cost is slight, because a fixed value chosen to fit the
frame period reaches the same image brightness that automatic exposure settles
on in good light; what is bought is a capture rate that is a property of the
configuration rather than of the hour of the day. The value is chosen per camera
and by eye, because two wrist cameras on one rig are not lit alike --- a single
figure that suits one over-exposes the other --- and because the judgement being
made is about a picture. The tool that sets them shows each stream's achieved
rate beside its image, so the one adjustment that costs rate is visible as it is
being made.

\paragraph{Recorded stream set.} The dataset records the colour views --- the
overhead, wrist and fingertip cameras --- the twelve-dimensional
proprioceptive state (five joints and a gripper per arm), the executed target
(Section~\ref{sec:contract}), and, when a depth-capable overhead camera is
fitted, its aligned depth image. That depth is sixteen-bit and is stored
losslessly outside the colour-video pipeline, because a depth image is not
three-channel colour and a colour-video encoder would corrupt it; the depth
scale and camera intrinsics are stored alongside. The full-rate side log
additionally holds the measured and target end-effector poses, the raw and
filtered operator input, and the gripper signals. Proprioception, command
and target share one unit convention and one channel layout, so a reader
never has to reconcile two coordinate conventions.

Each episode also carries an identifier taken from the wall clock at which it
began. Its position in the dataset cannot serve as its name, because removing
one episode renumbers every later one, so the same position denotes a different
recording afterwards, and any note an operator made against a number silently
comes to mean something else. The identifier is fixed when recording starts and
travels with the episode through renumbering, which lets a session log, a
review decision or a defect report keep pointing at the recording it was
written about.

An episode is also committed to storage the moment it is accepted, rather than
at the end of the session. This is not housekeeping: a dataset library that
accumulates episode records in memory and writes them in batches leaves the
session unreadable for as long as it runs, and loses every record still held in
memory if the session ends abnormally~--- and a recording whose record never
reached storage cannot be trained on, however complete its frames and video are
on the drive. A session lasting hours, on a rig where a camera can leave the bus
without warning, is precisely the case where that loss is likely and expensive.
Committing at episode granularity bounds any interruption to the recording in
flight, and has the further benefit that a session can be reviewed while it is
still being collected, so an operator discovers a framing or lighting mistake
during the sitting rather than after it.

\paragraph{Phase flag and gating.} Each frame carries a flag stating
whether teleoperation drove it. Homing moves and other non-teleoperation
motion are thereby marked and can be excluded when a policy is trained, so
the policy does not learn to freeze from frames where the recorded intent
fell back to the measured state (Section~\ref{sec:contract}). Gating then acts
at the episode level, and it is graded rather than absolute. A consumer bus
camera occasionally disappears for a second or two and returns by itself, and
discarding a long, otherwise good demonstration for such a blink costs an
operator far more than the missing second does. The recorder therefore holds a
stream loss: it stops adding frames, waits, and resumes the same episode once
every stream delivers again, abandoning the episode only if the stream stays
away. A session fault still discards outright.

We do not let that recovery hide a defect. A held episode really does contain a
gap, and because frame timestamps follow the frame index, the saved episode
would otherwise present one frame period across a hole that lasted seconds. The
recorder counts every hold and reports the total with the episode, so an
operator can inspect the episode and remove it if the motion jumps across the
join. Any end the operator did not ask for also sounds an audible cue, because
an operator attending to the arms cannot see the terminal, and would otherwise
continue demonstrating into a recording that had already stopped.
\unjustified{} the waiting period is set from the observed camera
reopen interval, not from any measurement of how large a gap a policy tolerates.

\paragraph{Session consistency.} A dataset is usually gathered over several
sittings, so its stream configuration must not drift between them. When a
session extends an existing dataset, the recorder reads back which streams
that dataset already holds --- the colour cameras, whether depth is present,
whether the end-effector channels were stored, and the frame rate --- and
follows them, rather than trusting flags typed afresh each time. The operator
therefore names the dataset and the task; the schema comes from the dataset
itself, and every episode in it shares one definition.
